\subsection*{Motion graphs}
\begin{enumerate}[label=\textbf{\thechapter.\arabic*},resume=exercises,leftmargin=*]
  \item An elevator's velocity--time graph rises steadily from $0$ to \SI{3.0}{\metre\per\second} over the first \SI{2.0}{\second}, stays at \SI{3.0}{\metre\per\second} for the next \SI{6.0}{\second}, then falls steadily to zero over the last \SI{3.0}{\second}. (a) Sketch the acceleration--time graph and give the acceleration in each phase. (b) Find the total displacement from the area under the velocity--time graph, and check it by integrating each piece. (c) Describe the shape of the position--time graph in each phase: straight or curved, and if curved, which way it bends.\label{ex:elevatorgraph}
  \item A cart's position--time graph is the parabola $x = 6t - t^2$ (metres, $t$ in seconds) for $0 \le t \le \SI{5}{\second}$. (a) When is the cart momentarily at rest, and where is it then? (b) Is the acceleration positive or negative? Find its value \elemtag\ from second differences of the positions at whole seconds and \calctag\ by differentiating twice. (c) Find the cart's displacement and its distance travelled over the five seconds. (d) Sketch the velocity--time and acceleration--time graphs.\label{ex:parabolagraph}
  \item $\star$ An object starts from rest at $x = 0$. Its acceleration--time graph is a step: $a_x = +\SI{2.0}{\metre\per\second\squared}$ from $t = 0$ to \SI{3.0}{\second}, then $a_x = \SI{-1.0}{\metre\per\second\squared}$ from \SI{3.0}{\second} to \SI{9.0}{\second}. (a) Sketch the velocity--time graph, finding the velocity at \SI{3.0}{\second} and at \SI{9.0}{\second} from areas under the acceleration--time graph. (b) When is the object farthest from its starting point? (c) Find its displacement at $t = \SI{9.0}{\second}$ \elemtag\ from the area of a triangle and \calctag\ by integrating $v_x(t)$ over each piece. (d) Sketch the position--time graph, describing the shape of each piece and what happens at the join.\label{ex:stepaccel}
  \item For each statement, say whether it is true, and if it is false give a counterexample drawn as a graph. (a) If the position--time graph is a straight line, the acceleration is zero. (b) If the velocity--time graph crosses the time axis, the object reverses direction at that instant. (c) If the position--time graph crosses the time axis, the velocity is zero at that instant. (d) The area under the acceleration--time graph is the displacement. (e) If the position--time graph is curved, the acceleration is nonzero. (f) The position--time graph of a ball thrown straight up is a picture of the ball's path.\label{ex:graphtruth}
\end{enumerate}

